\documentclass[10pt]{scrartcl}
\usepackage[mdthm,colorsec,noauthor,noasy]{evan}
\usepackage{handout}
\newcommand{\ip}[2]{\langle #1,#2\rangle}
\DeclareMathOperator{\tr}{tr}\DeclareMathOperator{\nul}{null}\DeclareMathOperator{\range}{range}\DeclareMathOperator{\spn}{span}
\newcommand{\LL}{\mathcal L}\newcommand{\PP}{\mathcal P}\newcommand{\MM}{\mathcal M}
\title{Math 110 --- Week 7}
\subtitle{Diagonalizability and the minimal polynomial (Notes Ch.~12--13, LADR 5C, 5D, 5B, HW7)}
\author{}\date{}\ihead{}
\begin{document}
\maketitle
\section{Facts used from earlier weeks}
$V$ is finite-dimensional over $\FF$ ($\RR$ or $\CC$), $n=\dim V$. Shorthand: ``$S$ kills $v$'' means $Sv=0$, and ``$q$ kills $T$'' means $q(T)=0$. From Week 6: eigenvalues, eigenvectors and eigenspaces $E(\lambda,T)=\nul(T-\lambda I)$, and invariant subspaces.

\begin{enumerate}[label=(T\arabic*),leftmargin=2.6em,itemsep=8pt]
  \ii \tool{Rank--nullity.} For linear $T\colon V\to W$: $\dim V=\dim\nul T+\dim\range T$.
  \why{\begin{steps}
    \ii Extend a basis $u_1,\dots,u_k$ of $\nul T$ to a basis $u_1,\dots,u_k,v_1,\dots,v_r$ of $V$. It suffices to show $Tv_1,\dots,Tv_r$ is a basis of $\range T$.
    \ii Spanning: $T$ of any vector is a combination of the $Tu_i$ (which are $0$) and the $Tv_j$.
    \ii Independent: if $\sum c_jTv_j=0$, then $\sum c_jv_j\in\nul T$ is a combination of the $u_i$, and independence of the basis of $V$ gives all $c_j=0$.
  \end{steps}}
  \ii \tool{For operators, injective $\iff$ surjective $\iff$ invertible.}
  \why{$\nul T=0\iff\dim\range T=n$ (T1) $\iff\range T=V$ (T3).}
  \ii \tool{Dimension shortcuts.}
  \begin{enumerate}[(a)]
    \ii An independent list of $n$ vectors is a basis.
    \ii A subspace of dimension $n$ is all of $V$.
    \ii If $U\cap W=\{0\}$ and $\dim U+\dim W=n$, then $V=U\oplus W$.
    \ii $\dim(U+W)=\dim U+\dim W-\dim(U\cap W)$.
  \end{enumerate}
  \why{\begin{steps}
    \ii (a), (b): an independent list extends to a basis, and every basis has length $n$.
    \ii (c): a basis of $U$ followed by a basis of $W$ is independent (a relation gives $u=-w\in U\cap W=\{0\}$) and has length $n$, so it is a basis by (a).
    \ii (d): extend a basis $x_1,\dots,x_k$ of $U\cap W$ to a basis $x_i,u_j$ of $U$ and, separately, to a basis $x_i,w_l$ of $W$. Together the $x_i,u_j,w_l$ span $U+W$; it remains to show they are independent.
    \ii If $x+u+w=0$ ($x,u,w$ combinations of the $x_i$, $u_j$, $w_l$), then $w=-x-u\in U\cap W$, so $w$ is a combination of the $x_i$. The basis $x_i,w_l$ of $W$ is independent, so $w=0$. Then $x+u=0$, and independence of $x_i,u_j$ gives $x=u=0$.
  \end{steps}}
  \ii \tool{Products of operators.} If $S,T$ are injective, so is $ST$. So if a product sends some $v\ne0$ to $0$, some factor is not injective.
  \why{$STv=0\Rightarrow Tv=0\Rightarrow v=0$.}
  \ii \tool{Polynomials in $T$ commute.} $p(T)q(T)=q(T)p(T)$. Hence $\nul p(T)$ and $\range p(T)$ are $T$-invariant.
  \why{Both products equal $(pq)(T)$. If $p(T)v=0$, then $p(T)(Tv)=T\,p(T)v=0$. And $T\bigl(p(T)u\bigr)=p(T)(Tu)\in\range p(T)$.}
  \ii \tool{Fundamental theorem of algebra.} A nonconstant $p\in\PP(\CC)$ factors as $c(z-\lambda_1)\cdots(z-\lambda_m)$ (cited). Over any field, $p(\lambda)=0\iff(z-\lambda)\mid p$.
  \why{Divide $p$ by $z-\lambda$: the remainder is a constant, and evaluating at $\lambda$ shows it equals $p(\lambda)$.}
  \ii \tool{Eigenvalues.}
  \begin{enumerate}[(a)]
    \ii $\lambda$ is an eigenvalue $\iff T-\lambda I$ is not injective.
    \ii If $Tv=\lambda v$, then $p(T)v=p(\lambda)v$.
    \ii Eigenvectors for distinct eigenvalues are independent (Prop.~54).
    \ii Over $\CC$ with $V\ne0$, every operator has an eigenvalue.
  \end{enumerate}
  \why{\begin{steps}
    \ii (a) is the definition: $Tv=\lambda v$ with $v\ne0$ means $(T-\lambda I)v=0$ with $v\ne0$.
    \ii (b): $T^kv=\lambda^kv$ by induction; add up the terms of $p$.
    \ii (d): take $v\ne0$. The $n+1$ vectors $v,Tv,\dots,T^nv$ are dependent, so $p(T)v=0$ for some $p\ne0$; $p$ is not constant since $v\ne0$.
    \ii Factor $p=c(z-\lambda_1)\cdots(z-\lambda_m)$ (T6). The product $c(T-\lambda_1I)\cdots(T-\lambda_mI)$ sends $v\ne0$ to $0$, so some $T-\lambda_jI$ is not injective (T4).
  \end{steps}}
  \ii \tool{Block diagonal matrices.} If $V=U\oplus W$ with both $T$-invariant, then in a basis of $U$ followed by one of $W$, $\MM(T)=\begin{psmallmatrix}A&0\\0&B\end{psmallmatrix}$, and $\nul T=\nul T|_U\oplus\nul T|_W$.
  \why{$Tu\in U$ is a combination of the $U$-basis only, so its column is $0$ in the $W$-rows (same for $W$). If $T(u+w)=0$, then $Tu=-Tw\in U\cap W=\{0\}$.}
  \ii \tool{Change of basis.} If column $j$ of $S$ holds the $B$-coordinates of the $j$-th vector of $B'$, then $\MM_{B'}(T)=S^{-1}\MM_B(T)S$. So similar matrices are exactly the matrices of one operator in two bases.
  \why{$S$ converts $B'$-coordinates to $B$-coordinates: apply $S$, then $\MM_B(T)$, then convert back with $S^{-1}$. Every invertible $S$ arises this way, since its columns form a basis.}
  \ii \tool{Quotient operator.} If $U$ is $T$-invariant, then $(T/U)(v+U)=Tv+U$ is a well-defined operator on $V/U$, $\dim V/U=\dim V-\dim U$, and $p(T/U)(v+U)=p(T)v+U$.
  \why{\begin{steps}
    \ii Well defined: if $v+U=v'+U$, then $v-v'\in U$, so $Tv-Tv'\in U$ ($U$ is invariant).
    \ii Dimension: $v\mapsto v+U$ maps $V$ onto $V/U$ with null space $U$; apply (T1).
    \ii $(T/U)^k(v+U)=T^kv+U$ by induction; add up the terms of $p$.
  \end{steps}}
\end{enumerate}

\section{This week}
\subsection{Upper-triangular matrices (Ch.~12.1)}

\begin{qbox}
  A diagonal matrix is the easiest kind to work with, but not every operator has one. What is the next best shape that every complex operator does have, and what can be read off from it?
  \begin{itemize}
    \ii $N=\begin{psmallmatrix}0&1\\0&0\end{psmallmatrix}$ has no diagonal matrix in any basis (Ex.~55 below), but it is upper-triangular, and its only eigenvalue $0$ is visible on the diagonal.
  \end{itemize}
\end{qbox}

\begin{definition*}[Def.~33]
  A square matrix $A=(a_{ij})$ is \vocab{upper-triangular} if $a_{ij}=0$ whenever $i>j$. Its \vocab{diagonal} is $a_{11},\dots,a_{nn}$. It is \vocab{diagonal} if every entry off the diagonal is $0$.
\end{definition*}
Ex.~53: $\begin{psmallmatrix}1&2&3&4\\0&5&6&7\\0&0&8&9\\0&0&0&10\end{psmallmatrix}$ is upper-triangular with diagonal $1,5,8,10$.

\begin{proposition*}[Prop.~59 / LADR 5.39]
  Let $v_1,\dots,v_n$ be a basis of $V$. The following are equivalent:
  \begin{enumerate}[(i)]
    \ii the matrix of $T$ in $v_1,\dots,v_n$ is upper-triangular;
    \ii $Tv_j\in\spn(v_1,\dots,v_j)$ for each $j$;
    \ii $\spn(v_1,\dots,v_j)$ is invariant under $T$ for each $j$.
  \end{enumerate}
\end{proposition*}
\begin{ideabox}
Column $j$ of the matrix lists the coordinates of $Tv_j$: $Tv_j=a_{1j}v_1+\dots+a_{nj}v_n$. ``Upper-triangular'' says the entries of column $j$ below row $j$ are $0$, which is the statement that $Tv_j$ uses only $v_1,\dots,v_j$. So the proposition only translates a statement about zeros in a matrix into a statement about where basis vectors are sent.

For $N=\begin{psmallmatrix}0&1\\0&0\end{psmallmatrix}$: $Ne_1=0$ and $Ne_2=e_1$. Each basis vector is sent to a combination of itself and earlier ones, never to a later one.
\end{ideabox}
\begin{proof}
  \begin{steps}
    \ii (i)$\Rightarrow$(ii): $Tv_j=a_{1j}v_1+\dots+a_{nj}v_n$, and $a_{ij}=0$ for $i>j$, so $Tv_j=a_{1j}v_1+\dots+a_{jj}v_j$.
    \ii (ii)$\Rightarrow$(iii): for $i\le j$, $Tv_i\in\spn(v_1,\dots,v_i)\subseteq\spn(v_1,\dots,v_j)$. So $T$ sends each of $v_1,\dots,v_j$, and hence their span, into $\spn(v_1,\dots,v_j)$.
    \ii (iii)$\Rightarrow$(i): $Tv_j\in\spn(v_1,\dots,v_j)$, so in $Tv_j=a_{1j}v_1+\dots+a_{nj}v_n$ the coefficients $a_{j+1,j},\dots,a_{nj}$ are $0$ (coordinates in a basis are unique).
  \end{steps}
\end{proof}

\begin{proposition*}[Prop.~60 / LADR 5.41]
  If $T$ has an upper-triangular matrix in some basis, the eigenvalues of $T$ are exactly the diagonal entries of that matrix.
\end{proposition*}
\begin{ideabox}
$\lambda$ is an eigenvalue exactly when $T-\lambda I$ is not injective (T7). The matrix of $T-\lambda I$ is again upper-triangular, with diagonal entries $\lambda_k-\lambda$. So the question is: when is an upper-triangular matrix injective?

Try $\begin{psmallmatrix}2&5\\0&3\end{psmallmatrix}\begin{psmallmatrix}x\\y\end{psmallmatrix}=0$. The bottom row says $3y=0$, so $y=0$; then the top row says $2x=0$, so $x=0$. Each row, read from the bottom up, forces one more coordinate to be $0$, as long as the diagonal entry in that row is not $0$. So a nonzero diagonal gives an injective matrix.

If instead some diagonal entry is $0$, say in position $j$, then $v_j$ is sent into $\spn(v_1,\dots,v_{j-1})$. So $j$ independent vectors $v_1,\dots,v_j$ are all sent into a space of dimension $j-1$, and that cannot be injective.
\end{ideabox}
\begin{proof}
  Let $\lambda_1,\dots,\lambda_n$ be the diagonal entries, in the basis $v_1,\dots,v_n$. Then $\MM(T-\lambda I)$ is upper-triangular with diagonal entries $\lambda_k-\lambda$.
  \begin{steps}
    \ii \emph{If $\lambda\ne\lambda_k$ for all $k$, then $T-\lambda I$ is injective.} Suppose $(T-\lambda I)v=0$ with $v=c_1v_1+\dots+c_nv_n$. Only $(T-\lambda I)v_n$ has a $v_n$-component, and it is $(\lambda_n-\lambda)v_n$. So the $v_n$-coefficient of $(T-\lambda I)v$ is $(\lambda_n-\lambda)c_n$, which must be $0$; hence $c_n=0$.
    \ii Now $v=c_1v_1+\dots+c_{n-1}v_{n-1}$, and the same argument with $v_{n-1}$ gives $c_{n-1}=0$. Continuing, every $c_k=0$. So $\lambda$ is not an eigenvalue.
    \ii \emph{If $\lambda=\lambda_j$, then $T-\lambda I$ is not injective.} For $i<j$, $(T-\lambda I)v_i\in\spn(v_1,\dots,v_i)$ (Prop.~59). And $(T-\lambda I)v_j\in\spn(v_1,\dots,v_{j-1})$, because its $v_j$-coefficient is $\lambda_j-\lambda=0$.
    \ii So $T-\lambda I$ maps the $j$-dimensional space $\spn(v_1,\dots,v_j)$ into the $(j-1)$-dimensional space $\spn(v_1,\dots,v_{j-1})$. By (T1) this restriction has a nonzero null space. So $\lambda$ is an eigenvalue.
  \end{steps}
\end{proof}
\begin{example*}[Ex.~54]
  Differentiation $D$ on $\PP_3(\RR)$ in the basis $1,x,x^2,x^3$ has matrix $\begin{psmallmatrix}0&1&0&0\\0&0&2&0\\0&0&0&3\\0&0&0&0\end{psmallmatrix}$, since $Dx^k=kx^{k-1}$. It is upper-triangular with diagonal $0,0,0,0$, so the only eigenvalue is $0$. Each basis vector is sent to a multiple of the previous one: $x^3\mapsto3x^2\mapsto6x\mapsto6\mapsto0$.
\end{example*}

\begin{theorem*}[Prop.~61 / LADR 5.47]
  If $V$ is a finite-dimensional complex vector space, every $T\in\LL(V)$ has an upper-triangular matrix in some basis.
\end{theorem*}
\begin{ideabox}
By Prop.~59 we need a basis in which each $Tv_k$ lies in the span of $v_k$ and the earlier vectors. Induction suggests: find a smaller invariant subspace $U$, take such a basis of $U$ (induction), and then add more vectors. The added vectors $v$ must satisfy $Tv\in U+\spn(v)$.

Which $U$ makes that automatic? Write $Tv=(T-\lambda I)v+\lambda v$. The second term is a multiple of $v$. The first term lies in $\range(T-\lambda I)$, whatever $v$ is. So take $U=\range(T-\lambda I)$: then \emph{every} vector $v$ satisfies $Tv\in U+\spn(v)$.

For induction $U$ has to be smaller than $V$, i.e.\ $T-\lambda I$ must not be surjective, which for operators is the same as not injective (T2). That is exactly the condition that $\lambda$ is an eigenvalue. Over $\CC$ an eigenvalue exists (T7); this is the only place $\CC$ is used.
\end{ideabox}
\begin{proof}
  Induction on $n=\dim V$; for $n=1$ every matrix is upper-triangular.
  \begin{steps}
    \ii Take an eigenvalue $\lambda$ of $T$ (T7) and let $U=\range(T-\lambda I)$. $U$ is invariant under $T$ (T5). $T-\lambda I$ is not injective, so it is not surjective (T2), so $\dim U<n$.
    \ii By induction there is a basis $u_1,\dots,u_m$ of $U$ in which $T|_U$ is upper-triangular: $Tu_i\in\spn(u_1,\dots,u_i)$.
    \ii Extend to a basis $u_1,\dots,u_m,v_1,\dots,v_k$ of $V$. Then $Tv_j=(T-\lambda I)v_j+\lambda v_j$, with $(T-\lambda I)v_j\in U=\spn(u_1,\dots,u_m)$. So $Tv_j\in\spn(u_1,\dots,u_m,v_1,\dots,v_j)$.
    \ii By Prop.~59 (ii)$\Rightarrow$(i), the matrix of $T$ in this basis is upper-triangular.
  \end{steps}
\end{proof}
\begin{itemize}
  \ii Over $\RR$ this can fail. In an upper-triangular basis, $Tv_1\in\spn(v_1)$, so $v_1$ is an eigenvector. The rotation $\begin{psmallmatrix}0&-1\\1&0\end{psmallmatrix}$ of $\RR^2$ has no eigenvector, so it has no upper-triangular matrix over $\RR$.
\end{itemize}

\subsection{Diagonalizable operators (Ch.~12.2)}

\begin{qbox}
  When is there a basis in which $T$ only rescales each basis vector?
  \begin{itemize}
    \ii $A=\begin{psmallmatrix}2&1\\1&2\end{psmallmatrix}$: $A(1,1)=3(1,1)$ and $A(1,-1)=(1,-1)$. In the basis $(1,1),(1,-1)$ the matrix is $\begin{psmallmatrix}3&0\\0&1\end{psmallmatrix}$.
    \ii $N=\begin{psmallmatrix}0&1\\0&0\end{psmallmatrix}$: every eigenvector is a multiple of $e_1$, so there are not enough eigenvectors for a basis.
  \end{itemize}
\end{qbox}

\begin{definition*}[Def.~34 / LADR 5.50]
  $T$ is \vocab{diagonalizable} if its matrix in some basis of $V$ is diagonal. A basis of $V$ consisting of eigenvectors of $T$ is an \vocab{eigenbasis}. The \vocab{geometric multiplicity} of $\lambda$ is $\dim E(\lambda,T)=\dim\nul(T-\lambda I)$; it is $0$ when $\lambda$ is not an eigenvalue.
\end{definition*}

\begin{theorem*}[Prop.~62 / LADR 5.55]
  Let $\lambda_1,\dots,\lambda_m$ be the distinct eigenvalues of $T$. The following are equivalent:
  \begin{enumerate}[(i)]
    \ii $T$ is diagonalizable;
    \ii $V$ has a basis of eigenvectors of $T$;
    \ii $V=E(\lambda_1,T)\oplus\dots\oplus E(\lambda_m,T)$;
    \ii $\dim V=\dim E(\lambda_1,T)+\dots+\dim E(\lambda_m,T)$.
  \end{enumerate}
\end{theorem*}
\begin{ideabox}
(i) and (ii) say the same thing: column $j$ of the matrix is the coordinates of $Tv_j$, and that column is $0$ except for the entry $\lambda$ in row $j$ exactly when $Tv_j=\lambda v_j$.

The content is in (iii) and (iv). The eigenspaces for different eigenvalues never overlap: their sum is always direct, because eigenvectors for distinct eigenvalues are independent (T7). So eigenspaces never give \emph{too many} vectors. The only way to fail is to give \emph{too few}, and that is a question of counting dimensions. That is why the single number in (iv) decides diagonalizability.

For $N=\begin{psmallmatrix}0&1\\0&0\end{psmallmatrix}$: the only eigenvalue is $0$ and $\dim E(0,N)=1<2$.
\end{ideabox}
\begin{proof}
  \begin{steps}
    \ii (i)$\iff$(ii): the matrix in $v_1,\dots,v_n$ is diagonal with entries $\mu_1,\dots,\mu_n$ exactly when $Tv_j=\mu_jv_j$ for every $j$, i.e.\ when every $v_j$ is an eigenvector.
    \ii \emph{The sum $E(\lambda_1,T)+\dots+E(\lambda_m,T)$ is always direct.} Suppose $u_1+\dots+u_m=0$ with $u_i\in E(\lambda_i,T)$. The nonzero $u_i$ are eigenvectors for distinct eigenvalues, so they are independent (T7); a sum of independent vectors is not $0$. So there are no nonzero $u_i$.
    \ii (ii)$\Rightarrow$(iii): each vector of an eigenbasis lies in some $E(\lambda_i,T)$, so the sum of the eigenspaces contains a basis of $V$ and equals $V$. It is direct by step 2.
    \ii (iii)$\Rightarrow$(iv): the dimension of a direct sum is the sum of the dimensions.
    \ii (iv)$\Rightarrow$(iii): by step 2 the sum of the eigenspaces is direct, so its dimension is $\sum\dim E(\lambda_i,T)=\dim V$. A subspace of full dimension is $V$ (T3).
    \ii (iii)$\Rightarrow$(ii): choose a basis of each $E(\lambda_i,T)$. Because the sum is direct, together they form a basis of $V$, and every vector in it is an eigenvector.
  \end{steps}
\end{proof}

\begin{example*}[Ex.~55: an operator that is not diagonalizable]
  Let $Tv_1=0$ and $Tv_k=v_{k-1}$ for $k\ge2$, with $n\ge2$. The matrix is upper-triangular with zero diagonal, so the only eigenvalue is $0$ (Prop.~60). $\nul T=\spn(v_1)$: if $T(\sum c_kv_k)=\sum_{k\ge2}c_kv_{k-1}=0$, then $c_2=\dots=c_n=0$. So $\sum_\lambda\dim E(\lambda,T)=1\ne n$, and $T$ is not diagonalizable by (iv).
\end{example*}

\begin{proposition*}[Prop.~63 / LADR 5.58]
  If $T$ has $n=\dim V$ distinct eigenvalues, then $T$ is diagonalizable.
\end{proposition*}
\begin{proof}
  Choose an eigenvector $v_i$ for each eigenvalue $\lambda_i$. They are independent (T7), and there are $n$ of them, so they form a basis (T3). It is an eigenbasis.
\end{proof}
\begin{itemize}
  \ii The condition is sufficient, not necessary: $I$ has one eigenvalue and is diagonal.
  \ii Ex.~56: diagonalizability depends on the field. $\begin{psmallmatrix}1&2\\2&3\end{psmallmatrix}$ has eigenvalues $2\pm\sqrt5$ (the roots of $x^2-4x-1$), so it is diagonalizable over $\RR$; over $\mathbb Q$ it has no eigenvalues. $\begin{psmallmatrix}0&-1\\1&0\end{psmallmatrix}$ has eigenvalues $\pm i$, so it is diagonalizable over $\CC$ but not over $\RR$.
\end{itemize}

\method{Deciding diagonalizability and diagonalizing a matrix $A$.}
\begin{steps}
  \ii Find the eigenvalues $\lambda_i$ (for a triangular matrix, the diagonal entries).
  \ii For each, $\dim E(\lambda_i,A)=n-\operatorname{rank}(A-\lambda_iI)$ (T1). $A$ is diagonalizable exactly when these add up to $n$ (Prop.~62).
  \ii If so, find a basis of each $\nul(A-\lambda_iI)$ and put all these vectors as the columns of $S$. Then $S^{-1}AS=D$ is diagonal, with the eigenvalues in the same order as the columns (T9).
\end{steps}
\begin{example*}
  $A=\begin{psmallmatrix}2&1\\1&2\end{psmallmatrix}$. $A-3I=\begin{psmallmatrix}-1&1\\1&-1\end{psmallmatrix}$ has null space $\spn((1,1))$; $A-I=\begin{psmallmatrix}1&1\\1&1\end{psmallmatrix}$ has null space $\spn((1,-1))$. Dimensions $1+1=2$. With $S=\begin{psmallmatrix}1&1\\1&-1\end{psmallmatrix}$: $S^{-1}AS=\begin{psmallmatrix}3&0\\0&1\end{psmallmatrix}$, and $A^k=S\begin{psmallmatrix}3^k&0\\0&1\end{psmallmatrix}S^{-1}$.
\end{example*}

\subsection{Iteration and eigenmodes (Ch.~12.3)}

\begin{proposition*}[Prop.~64: iteration in an eigenbasis]
  If $v_1,\dots,v_n$ is an eigenbasis with $Tv_j=\lambda_jv_j$ and $v=a_1v_1+\dots+a_nv_n$, then $T^kv=a_1\lambda_1^kv_1+\dots+a_n\lambda_n^kv_n$ for every $k\ge0$.
\end{proposition*}
\begin{proof}
  $T^kv_j=\lambda_j^kv_j$ by induction on $k$; apply $T^k$ to $v$ term by term.
\end{proof}
The notes call each $a_jv_j$ an \vocab{eigenmode}. Under repeated application of $T$, a mode with $|\lambda_j|<1$ shrinks to $0$, one with $|\lambda_j|=1$ keeps its size, and one with $|\lambda_j|>1$ grows.

\begin{example*}[Markov chain]
  $P=\begin{psmallmatrix}1/2&1/4&0\\1/2&1/2&1/2\\0&1/4&1/2\end{psmallmatrix}$, $\pi=\begin{psmallmatrix}1/4\\1/2\\1/4\end{psmallmatrix}$, $u=\begin{psmallmatrix}-1\\0\\1\end{psmallmatrix}$, $w=\begin{psmallmatrix}1\\-2\\1\end{psmallmatrix}$.
  \begin{steps}
    \ii Multiplying: $P\pi=\pi$, $Pu=\frac12u$, $Pw=0$. So $\pi,u,w$ is an eigenbasis with eigenvalues $1,\frac12,0$.
    \ii The start $p_0=(1,0,0)$ is $\pi-\frac12u+\frac14w$ (check the first coordinate: $\frac14+\frac12+\frac14=1$).
    \ii By Prop.~64, for $k\ge1$: $P^kp_0=\pi-\frac12\cdot\frac1{2^k}u+0=\pi-2^{-(k+1)}u$.
    \ii So $P^kp_0\to\pi$, and the distance to $\pi$ halves at every step. The eigenvalue $1$ gives the limit; the next largest eigenvalue, $\frac12$, gives the speed.
  \end{steps}
\end{example*}
\begin{itemize}
  \ii \textbf{Spectral gap.} $P_\varepsilon=\begin{psmallmatrix}1-\varepsilon&\varepsilon\\\varepsilon&1-\varepsilon\end{psmallmatrix}$, $0<\varepsilon<\frac12$, has eigenvectors $(1,1)$ and $(1,-1)$ with eigenvalues $1$ and $1-2\varepsilon$. So $P_\varepsilon^k(1,-1)=(1-2\varepsilon)^k(1,-1)$. When $\varepsilon$ is small (moving between the two states is rare), $1-2\varepsilon$ is close to $1$ and the starting state is forgotten slowly. The number $1-(1-2\varepsilon)=2\varepsilon$ is the \vocab{spectral gap}.
  \ii \textbf{PageRank} (Week 6): $\RR^n=\spn(r)\oplus H$ with $Gr=r$ and $Gx=\alpha Px$ for $x\in H$. If $x\in H$ is an eigenvector of $P$ with eigenvalue $\lambda$, then $Gx=\alpha\lambda x$. Every eigenvalue of $P$ has modulus $\le1$, so these eigenvalues of $G$ have modulus $\le\alpha$: every mode other than $r$ shrinks by at least the factor $\alpha$ per step.
\end{itemize}
\begin{example*}[Fibonacci numbers, LADR 5D Exercise 21]
  $F_0=0$, $F_1=1$, $F_{k+1}=F_k+F_{k-1}$. Let $T(x,y)=(y,x+y)$, so $T(F_{k-1},F_k)=(F_k,F_{k+1})$ and $T^k(0,1)=(F_k,F_{k+1})$.
  \begin{steps}
    \ii $T(1,\lambda)=(\lambda,1+\lambda)$ equals $\lambda(1,\lambda)$ exactly when $\lambda^2=\lambda+1$: $\varphi=\frac{1+\sqrt5}2$ and $\psi=\frac{1-\sqrt5}2$. Two distinct eigenvalues, so $(1,\varphi),(1,\psi)$ is an eigenbasis (Prop.~63).
    \ii $(0,1)=\frac1{\sqrt5}\bigl((1,\varphi)-(1,\psi)\bigr)$, since $\varphi-\psi=\sqrt5$.
    \ii By Prop.~64, $T^k(0,1)=\frac1{\sqrt5}\bigl(\varphi^k(1,\varphi)-\psi^k(1,\psi)\bigr)$. First coordinate: $F_k=\dfrac{\varphi^k-\psi^k}{\sqrt5}$.
  \end{steps}
\end{example*}

\begin{shortcuts}
  \ii \textbf{Triangular matrix}: the eigenvalues are the diagonal entries (Prop.~60). No computation.
  \ii \textbf{$n$ distinct eigenvalues}: diagonalizable (Prop.~63). A triangular matrix with distinct diagonal entries is diagonalizable.
  \ii \textbf{One eigenvalue only}: then $T$ is diagonalizable exactly when $T=\lambda I$ (a diagonal matrix with every diagonal entry $\lambda$ is $\lambda I$). So a triangular matrix with constant diagonal that is not $\lambda I$ is not diagonalizable.
  \ii \textbf{Dimension of an eigenspace by rank}: $\dim E(\lambda,A)=n-\operatorname{rank}(A-\lambda I)$.
  \ii \textbf{Powers}: write $v$ in an eigenbasis and multiply each coefficient by $\lambda_j^k$ (Prop.~64); or $A^k=SD^kS^{-1}$.
\end{shortcuts}
\begin{trapbox}
  \ii Repeated diagonal entries decide nothing by themselves: $I$ and $\begin{psmallmatrix}1&1\\0&1\end{psmallmatrix}$ both have diagonal $1,1$; the first is diagonalizable, the second is not.
  \ii Upper-triangular is not the same as diagonalizable. Every complex operator has an upper-triangular matrix; not every one has a diagonal matrix.
  \ii Over $\RR$ an operator may have no upper-triangular matrix at all (rotation).
  \ii Diagonalizability depends on the field (Ex.~56).
  \ii The eigenvalues are the diagonal entries only for a \emph{triangular} matrix: $\begin{psmallmatrix}0&1\\1&0\end{psmallmatrix}$ has diagonal $0,0$ but eigenvalues $\pm1$.
\end{trapbox}



\subsection{Minimal polynomial (Ch.~13)}

\begin{qbox}
  Powers of $T$ cannot stay independent forever, so some nonzero polynomial kills $T$. The one of smallest degree tells you the eigenvalues of $T$ and whether $T$ is diagonalizable.
  \begin{itemize}
    \ii $T=\begin{psmallmatrix}2&0\\0&3\end{psmallmatrix}$: $(T-2I)(T-3I)=0$, and no polynomial of degree $1$ kills $T$, since $T$ is not a multiple of $I$. The roots $2,3$ are the eigenvalues.
    \ii $N=\begin{psmallmatrix}0&1\\0&0\end{psmallmatrix}$: $N^2=0$ but $N\ne0$, so the smallest polynomial is $x^2$. Its repeated root is what stops $N$ from being diagonalizable.
  \end{itemize}
\end{qbox}

\begin{definition*}[Def.~35, 36 / LADR 5.14, 5.22, 5.24]
  For $p(x)=a_mx^m+\dots+a_1x+a_0$, $p(T)=a_mT^m+\dots+a_1T+a_0I$. Let $k$ be the smallest index with $T^k\in\spn(I,T,\dots,T^{k-1})$, and write $T^k=a_{k-1}T^{k-1}+\dots+a_0I$. The \vocab{minimal polynomial} of $T$ is $m_T(x)=x^k-a_{k-1}x^{k-1}-\dots-a_0$. (The notes call it $p$.)
\end{definition*}
\begin{itemize}
  \ii $k$ exists: $\LL(V)$ has dimension $n^2$ (operators are $n\times n$ matrices), so the $n^2+1$ operators $I,T,\dots,T^{n^2}$ are dependent. Solving a dependence for the highest power with a nonzero coefficient puts that power in the span of the lower ones.
  \ii $I,T,\dots,T^{k-1}$ are independent, by the same argument: a dependence among them would put some $T^j$ with $j<k$ in the span of lower powers. So the $a_i$ are unique.
  \ii Shorthand used below: $q$ \emph{kills} $T$ means $q(T)=0$; an operator $S$ \emph{kills} $v$ means $Sv=0$.
  \ii $m_T$ is monic and $m_T(T)=0$. No nonzero polynomial of degree $<k$ kills $T$, since it would be a dependence among $I,\dots,T^{k-1}$. So $m_T$ is the unique monic polynomial of least degree that kills $T$ (two of them would differ by a nonzero polynomial of lower degree that kills $T$).
  \ii Ex.~57: $m_I=x-1$. For $T=\begin{psmallmatrix}2&1\\1&2\end{psmallmatrix}$, $I,T$ are independent and $T^2=\begin{psmallmatrix}5&4\\4&5\end{psmallmatrix}=4T-3I$, so $m_T=x^2-4x+3=(x-1)(x-3)$.
\end{itemize}

\begin{proposition*}[Prop.~65 / LADR 5.29]
  For $q\in\PP(\FF)$: $q(T)=0\iff m_T\mid q$.
\end{proposition*}
\begin{ideabox}
To show: Every polynomial that kills $T$ is a multiple of $m_T$.

$m_T$ has the smallest degree among nonzero polynomials that kill $T$. Polynomials can be divided with remainder, and the remainder has smaller degree than the divisor.

Divide $q$ by $m_T$. The remainder kills $T$ too, and its degree is too small for it to be nonzero.
\end{ideabox}
\begin{proof}
  \begin{steps}
    \ii If $q=m_Tr$, then $q(T)=m_T(T)r(T)=0$ (T5).
    \ii If $q(T)=0$, divide: $q=m_Tr+s$ with $s=0$ or $\deg s<\deg m_T$. Then $s(T)=q(T)-m_T(T)r(T)=0$.
    \ii A nonzero $s$ of degree $<\deg m_T$ cannot kill $T$. So $s=0$ and $m_T\mid q$.
  \end{steps}
\end{proof}

\begin{proposition*}[Prop.~66 / LADR 5.27]
  The eigenvalues of $T$ are exactly the roots of $m_T$ in $\FF$. If $m_T=(x-\lambda)s$, then $s(T)\ne0$, and every nonzero vector $s(T)w$ is an eigenvector for $\lambda$.
\end{proposition*}
\begin{ideabox}
To show: $m_T(\lambda)=0$ for each eigenvalue $\lambda$.

On an eigenvector, a polynomial in $T$ acts as the number $p(\lambda)$ (T7). Apply this to $m_T$, which is $0$ as an operator.

To show: An eigenvector for a root $\lambda$, i.e.\ a nonzero vector killed by $T-\lambda I$.

Factor $m_T=(x-\lambda)s$. Then $(T-\lambda I)s(T)=m_T(T)=0$, so $T-\lambda I$ kills every vector of the form $s(T)w$. Such a vector can be chosen nonzero because $s$ has smaller degree than $m_T$, so $s(T)\ne0$.
\end{ideabox}
\begin{proof}
  \begin{steps}
    \ii If $Tv=\lambda v$ with $v\ne0$, then $0=m_T(T)v=m_T(\lambda)v$, so $m_T(\lambda)=0$.
    \ii If $m_T(\lambda)=0$, write $m_T=(x-\lambda)s$ (T6). Since $\deg s<\deg m_T$, $s(T)\ne0$; pick $w$ with $s(T)w\ne0$.
    \ii $(T-\lambda I)\bigl(s(T)w\bigr)=m_T(T)w=0$, so $s(T)w$ is an eigenvector for $\lambda$.
  \end{steps}
  The notes argue the second half by contradiction: if $T-\lambda I$ were invertible, multiplying $0=s(T)(T-\lambda I)$ by $(T-\lambda I)^{-1}$ would give $s(T)=0$.
\end{proof}

\begin{proposition*}[Prop.~67 / LADR 5.22]
  $\deg m_T\le n$.
\end{proposition*}
\begin{ideabox}
To show: Some polynomial of degree $\le n$ that kills $T$. Then $m_T$ divides it (Prop.~65).

For one vector $v\ne0$, the $n+1$ vectors $v,Tv,\dots,T^nv$ are dependent, so $q(T)v=0$ for some $q\ne0$ of degree $\le n$. This kills $v$, but not necessarily all of $V$.

The obvious first attempt: Do this for each basis vector and multiply the $q$'s. The product kills $T$, but its degree can be about $n^2$.

Induction on $\dim V$. $U=\nul q(T)$ is invariant and contains $v$. If $U\ne V$, both $U$ and $V/U$ are smaller than $V$. By induction, a polynomial $r$ of degree $\le\dim U$ kills $T$ on $U$, and a polynomial $p$ of degree $\le\dim V/U$ kills $T/U$, which means $p(T)$ sends $V$ into $U$. So $r(T)p(T)$ kills $V$, and its degree is $\le\dim U+\dim V/U=n$.
\end{ideabox}
\begin{proof}
  \begin{steps}
    \ii If $\dim V=1$, $T=cI$ and $m_T=x-c$. In general, fix $v\ne0$ and $q\ne0$ with $\deg q\le n$ and $q(T)v=0$. $U=\nul q(T)$ is invariant (T5).
    \ii If $q(T)=0$, then $m_T\mid q$, so $\deg m_T\le n$.
    \ii Otherwise $0<\dim U<n$. By induction, $r=m_{T|_U}$ has $\deg r\le\dim U$ and $p=m_{T/U}$ has $\deg p\le\dim V/U$ (T10).
    \ii For $v\in V$: $p(T)v+U=p(T/U)(v+U)=0+U$, so $p(T)v\in U$, so $r(T)p(T)v=0$. Hence $m_T\mid rp$, and $\deg(rp)\le\dim U+\dim V/U=n$ (T10).
  \end{steps}
\end{proof}
\needspace{7\baselineskip}\textbf{Second proof (LADR 5.22), without quotient spaces.} Same plan, with $\range q(T)$ in place of $V/U$. Induction on $n$.
\begin{steps}
  \ii Fix $v\ne0$ and let $q$ be the monic polynomial of least degree $d$ with $q(T)v=0$; $d\ge1$ since $v\ne0$.
  \ii $v,Tv,\dots,T^{d-1}v$ are independent (a dependence would be a smaller polynomial killing $v$), and they lie in $\nul q(T)$, because $q(T)T^jv=T^jq(T)v=0$. So $\dim\nul q(T)\ge d$.
  \ii $W=\range q(T)$ is invariant (T5) and $\dim W\le n-d<n$ (T1). By induction some $s$ of degree $\le n-d$ kills $T|_W$.
  \ii For every $u$: $q(T)u\in W$, so $s(T)q(T)u=0$. So $sq$ kills $T$, and $\deg(sq)\le n$.
\end{steps}
\begin{itemize}
  \ii Ex.~58: every degree from $1$ to $n$ occurs.
  \begin{itemize}
    \ii $m_I=x-1$ in every dimension.
    \ii The shift $N$ ($Ne_1=0$, $Ne_j=e_{j-1}$) has $m_N=x^n$, since $N^{n-1}e_n=e_1\ne0$ and $N^n=0$.
    \ii Degree $k$, $1\le k\le n$: let $Ne_1=0$, $Ne_j=e_{j-1}$ for $2\le j\le k$, $Ne_j=0$ for $j>k$. Then $N^k=0$ and $N^{k-1}e_k=e_1\ne0$, so $m_N=x^k$.
  \end{itemize}
\end{itemize}

\begin{proposition*}[Prop.~68 / LADR 5.62]
  $T$ is diagonalizable $\iff m_T=(x-\lambda_1)\cdots(x-\lambda_m)$ with $\lambda_1,\dots,\lambda_m\in\FF$ distinct.
\end{proposition*}
\begin{ideabox}
To show: $\prod_i(T-\lambda_iI)=0$, where $\lambda_1,\dots,\lambda_m$ are the distinct eigenvalues.

On an eigenvector for $\lambda_j$, the factor $T-\lambda_jI$ gives $0$, and the factors commute, so let that factor act first. There is an eigenbasis, so the product kills $V$. Then $m_T$ divides $\prod(x-\lambda_i)$ and contains every $x-\lambda_i$ (Prop.~66), so it equals it.

To show: Write every $v$ as a sum of eigenvectors.

$(T-\lambda_jI)\prod_{i\ne j}(T-\lambda_iI)=m_T(T)=0$. So every vector of the form $\prod_{i\ne j}(T-\lambda_iI)w$ is in $E(\lambda_j,T)$. This is how to produce eigenvectors.

Scale these so that they add up to $v$. The polynomial $g_j(x)=\prod_{i\ne j}\frac{x-\lambda_i}{\lambda_j-\lambda_i}$ is $1$ at $\lambda_j$ and $0$ at the other $\lambda_i$. Then $\sum_jg_j-1$ has degree $<m$ and $m$ roots, so $\sum_jg_j=1$. Apply both sides to $v$.
\end{ideabox}
\begin{proof}
  \begin{steps}
    \ii ($\Rightarrow$) Let $v_1,\dots,v_n$ be an eigenbasis, $Tv_k=\mu v_k$ with $\mu$ one of the $\lambda_i$. Letting the factor $T-\mu I$ act first, $\prod_i(T-\lambda_iI)v_k=0$ for every $k$ (T5). So $m_T\mid\prod_i(x-\lambda_i)$ (Prop.~65). Each $\lambda_i$ is a root of $m_T$ (Prop.~66), so $m_T=\prod_i(x-\lambda_i)$.
    \ii ($\Leftarrow$) Define $g_j$ as above. $\sum_jg_j-1$ has degree $\le m-1$ and vanishes at the $m$ points $\lambda_1,\dots,\lambda_m$. A nonzero polynomial of degree $\le m-1$ has at most $m-1$ roots (each root splits off a factor, T6), so $\sum_jg_j=1$.
    \ii For $v\in V$: $v=\sum_jg_j(T)v$, and $(T-\lambda_jI)g_j(T)v=c_j\,m_T(T)v=0$, where $c_j=\prod_{i\ne j}(\lambda_j-\lambda_i)^{-1}$. So $g_j(T)v\in E(\lambda_j,T)$.
    \ii So $V=E(\lambda_1,T)+\dots+E(\lambda_m,T)$. This sum is direct (T7), so bases of the eigenspaces together form an eigenbasis.
  \end{steps}
\end{proof}
\begin{fig}
\begin{tikzpicture}
\begin{axis}[width=8.6cm,height=5.2cm,xmin=0.6,xmax=3.4,ymin=-1.1,ymax=1.8,axis lines=left,
  xtick={1,2,3},xticklabels={$\lambda_1=1$,$\lambda_2=2$,$\lambda_3=3$},ytick={0,1},
  domain=0.6:3.4,samples=80,legend style={font=\small,draw=none,at={(1.03,1)},anchor=north west},legend cell align=left]
  \addplot[blue,thick] {(x-2)*(x-3)/2}; \addlegendentry{$g_1$}
  \addplot[red,thick] {-(x-1)*(x-3)}; \addlegendentry{$g_2$}
  \addplot[OliveGreen,thick] {(x-1)*(x-2)/2}; \addlegendentry{$g_3$}
  \addplot[black,dashed,thick] {1}; \addlegendentry{$g_1+g_2+g_3$}
  \addplot[gray,thin] {0};
  \addplot[only marks,mark=*,mark size=1.6pt,black] coordinates {(1,1) (2,1) (3,1) (1,0) (2,0) (3,0)};
\end{axis}
\end{tikzpicture}
\figcap{$g_j$ is $1$ at $\lambda_j$ and $0$ at the other eigenvalues, and $g_1+g_2+g_3=1$. Applying $1=\sum g_j$ to $T$ writes $v=\sum g_j(T)v$, and each $g_j(T)v$ lies in $E(\lambda_j,T)$ (Prop.~68). The same $g_j(T)$ are the projections of F9.}
\end{fig}
\needspace{7\baselineskip}\textbf{Second proof of ($\Leftarrow$), the one in the notes.} Instead of splitting $v$ into all its eigenvector pieces at once, split off the $\lambda_1$-piece and use induction for the rest. Let $g=\prod_{i\ge2}(x-\lambda_i)$, so $m_T=(x-\lambda_1)g$ and $g(\lambda_1)\ne0$.
\begin{steps}
  \ii Dividing by $x-\lambda_1$ leaves the remainder $g(\lambda_1)$ (T6): $g=(x-\lambda_1)q+g(\lambda_1)$.
  \ii For $v\in V$ let $u=g(T)v/g(\lambda_1)$. Then $(T-\lambda_1I)u=m_T(T)v/g(\lambda_1)=0$, so $u\in\nul(T-\lambda_1I)$, and $v-u=-(T-\lambda_1I)q(T)v/g(\lambda_1)\in\range(T-\lambda_1I)$.
  \ii So $V=\nul(T-\lambda_1I)+\range(T-\lambda_1I)$. The two dimensions add up to $n$ (T1), so the intersection is $0$ (last item of T3) and the sum is direct.
  \ii $\range(T-\lambda_1I)$ is invariant, smaller than $V$ (since $\nul(T-\lambda_1I)\ne0$: $\lambda_1$ is an eigenvalue by Prop.~66), and its minimal polynomial divides $m_T$ (F3), so it again has distinct linear factors. By induction $T$ restricted to it has an eigenbasis; add a basis of $E(\lambda_1,T)$.
\end{steps}


\subsubsection*{More facts about $m_T$}
Throughout, write $m_T=x^d+c_{d-1}x^{d-1}+\dots+c_1x+c_0$.

\begin{proposition*}[F1: invertibility, HW7 P3]
  $T$ is invertible $\iff c_0\ne0$. In that case $T^{-1}=-\dfrac1{c_0}\bigl(T^{d-1}+c_{d-1}T^{d-2}+\dots+c_1I\bigr)$, a polynomial in $T$.
\end{proposition*}
\begin{ideabox}
To show: $T^{-1}$ written with powers of $T$.

$m_T(T)=0$, and every term of $m_T(T)$ except $c_0I$ contains a factor $T$.

Move $c_0I$ to the other side and factor out $T$. That gives $T\cdot(\text{something})=-c_0I$, so $\frac{-1}{c_0}(\text{something})$ is the inverse.
\end{ideabox}
\begin{proof}
  \begin{steps}
    \ii $c_0=m_T(0)$. By Prop.~66, $c_0=0\iff0$ is an eigenvalue $\iff T$ is not injective (T7) $\iff T$ is not invertible (T2).
    \ii If $c_0\ne0$: $m_T(T)=0$ says $T^d+c_{d-1}T^{d-1}+\dots+c_1T=-c_0I$.
    \ii Factor $T$ out of the left side: $T\bigl(T^{d-1}+c_{d-1}T^{d-2}+\dots+c_1I\bigr)=-c_0I$. Divide by $-c_0$.
  \end{steps}
\end{proof}

\begin{proposition*}[F2: degree one, M4]
  If $V\ne\{0\}$: $\deg m_T=1\iff T=cI$ for some $c\in\FF$.
\end{proposition*}
\begin{proof}
  \begin{steps}
    \ii ($\Rightarrow$) $m_T=x-c$ and $m_T(T)=0$ say $T-cI=0$.
    \ii ($\Leftarrow$) $x-c$ kills $cI$. No polynomial of degree $0$ does: a nonzero constant $a$ gives $aI\ne0$ since $V\ne\{0\}$. So $m_{cI}=x-c$.
  \end{steps}
\end{proof}

\begin{proposition*}[F3: restrictions and quotients, Ex.~60]
  Let $U$ be $T$-invariant.
  \begin{enumerate}[(a)]
    \ii $m_{T|_U}\mid m_T$ and $m_{T/U}\mid m_T$.
    \ii A monic divisor of a product $(x-\lambda_1)\cdots(x-\lambda_m)$ of distinct linear factors is the product of some of those factors.
    \ii So if $T$ is diagonalizable, so are $T|_U$ and $T/U$.
  \end{enumerate}
\end{proposition*}
\begin{proof}
  (a)
  \begin{steps}
    \ii For $u\in U$: $m_T(T|_U)u=m_T(T)u=0$. So $m_T$ kills $T|_U$, and $m_{T|_U}\mid m_T$ by Prop.~65 applied to $T|_U$.
    \ii For $v+U\in V/U$: $m_T(T/U)(v+U)=m_T(T)v+U=0+U$ (T10). So $m_{T/U}\mid m_T$.
  \end{steps}
  (b) Let $s$ be the monic divisor, $st=\prod_{i=1}^m(x-\lambda_i)$.
  \begin{steps}
    \ii Each $\lambda_i$ makes the product $0$, so it is a root of $s$ or of $t$. Let $S$ be the set of $\lambda_i$ that are roots of $s$, and $R$ the remaining $\lambda_i$ (roots of $t$).
    \ii Distinct roots split off one factor each (T6), so $\prod_{\lambda\in S}(x-\lambda)$ divides $s$ and $\prod_{\lambda\in R}(x-\lambda)$ divides $t$. Hence $\deg s\ge|S|$ and $\deg t\ge|R|$.
    \ii But $\deg s+\deg t=m=|S|+|R|$, so both are equalities. A monic $s$ divisible by a monic polynomial of the same degree equals it: $s=\prod_{\lambda\in S}(x-\lambda)$.
  \end{steps}
  (c) $m_T$ is a product of distinct linear factors (Prop.~68). By (a) and (b), so are $m_{T|_U}$ and $m_{T/U}$. Apply Prop.~68 again.
\end{proof}

\begin{proposition*}[F4: direct sums, block diagonal matrices]
  If $V=U_1\oplus\dots\oplus U_r$ with every $U_i$ invariant, then $m_T=\operatorname{lcm}\bigl(m_{T|_{U_1}},\dots,m_{T|_{U_r}}\bigr)$. For a block diagonal matrix, $m$ is the lcm of the minimal polynomials of the blocks.
\end{proposition*}
\begin{proof}
  \begin{steps}
    \ii Write $v=u_1+\dots+u_r$ with $u_i\in U_i$. Then $q(T)v=\sum_iq(T|_{U_i})u_i$, so $q(T)=0\iff q(T|_{U_i})=0$ for every $i$.
    \ii By Prop.~65, that holds $\iff$ every $m_{T|_{U_i}}$ divides $q$.
    \ii So $m_T$ is the monic polynomial of least degree divisible by all the $m_{T|_{U_i}}$: their lcm. The blocks of the matrix are the matrices of the $T|_{U_i}$ (T8).
  \end{steps}
\end{proof}

\begin{proposition*}[F5: similarity, transpose, dual]
  Similar matrices have the same minimal polynomial. Also $m_{A^{\mathsf T}}=m_A$ and $m_{T'}=m_T$.
\end{proposition*}
\begin{proof}
  \begin{steps}
    \ii $(S^{-1}AS)^k=S^{-1}A^kS$ (the inner factors $SS^{-1}$ cancel), so $p(S^{-1}AS)=S^{-1}p(A)S$. This is $0$ exactly when $p(A)=0$. So the same polynomials kill both, and the minimal polynomials agree.
    \ii $(A^{\mathsf T})^k=(A^k)^{\mathsf T}$, so $p(A^{\mathsf T})=p(A)^{\mathsf T}$, which is $0$ exactly when $p(A)=0$.
    \ii In dual bases, $\MM(T')=\MM(T)^{\mathsf T}$; apply step 2.
  \end{steps}
\end{proof}

\begin{proposition*}[F6: shifts, scalings, inverses, HW7 P4, M5]
  With $d=\deg m_T$:
  \begin{enumerate}[(a)]
    \ii $m_{T-\lambda I}(x)=m_T(x+\lambda)$;
    \ii $m_{cT}(x)=c^d\,m_T(x/c)$ for $c\ne0$;
    \ii if $T$ is invertible, $m_{T^{-1}}(x)=\frac1{c_0}x^dm_T(1/x)$, i.e.\ reverse the list of coefficients of $m_T$ and divide by $c_0$.
  \end{enumerate}
\end{proposition*}
\begin{ideabox}
To show: The minimal polynomial of a new operator built from $T$.

A polynomial in the new operator is a polynomial in $T$ after a substitution: $q(T-\lambda I)$ is $r(T)$ with $r(x)=q(x-\lambda)$.

The substitution can be undone and does not change degrees. So it matches the polynomials that kill the new operator with the polynomials that kill $T$, and the one of least degree goes to the one of least degree.
\end{ideabox}
\begin{proof}
  (a)
  \begin{steps}
    \ii For any $q$, $q(T-\lambda I)=r(T)$ with $r(x)=q(x-\lambda)$.
    \ii So $q$ kills $T-\lambda I$ $\iff m_T(x)\mid q(x-\lambda)$ (Prop.~65) $\iff m_T(x+\lambda)\mid q(x)$ (replace $x$ by $x+\lambda$ throughout).
    \ii $m_T(x+\lambda)$ is monic, so it is $m_{T-\lambda I}$.
  \end{steps}
  (b) The same three steps with $r(x)=q(cx)$: $q$ kills $cT\iff m_T(x/c)\mid q(x)$. Multiply $m_T(x/c)$ by $c^d$ to make it monic.

  (c)
  \begin{steps}
    \ii Let $r(x)=\frac1{c_0}x^dm_T(1/x)=\frac1{c_0}\bigl(c_0x^d+c_1x^{d-1}+\dots+c_{d-1}x+1\bigr)$. It is monic of degree $d$ ($c_0\ne0$ by F1).
    \ii $r(T^{-1})=\frac1{c_0}\bigl(c_0T^{-d}+c_1T^{-(d-1)}+\dots+I\bigr)=\frac1{c_0}T^{-d}m_T(T)=0$. So $m_{T^{-1}}\mid r$ and $\deg m_{T^{-1}}\le d$.
    \ii The same argument with $T^{-1}$ in place of $T$ (its inverse is $T$) gives $d\le\deg m_{T^{-1}}$.
    \ii So $m_{T^{-1}}$ divides $r$ and has the same degree; both are monic, so they are equal.
  \end{steps}
\end{proof}

\begin{proposition*}[F7: one vector at a time]
  For $v\in V$, let $p_v$ be the monic polynomial of least degree with $p_v(T)v=0$.
  \begin{enumerate}[(a)]
    \ii $q(T)v=0\iff p_v\mid q$. In particular $p_v\mid m_T$.
    \ii For a basis $v_1,\dots,v_n$: $m_T=\operatorname{lcm}(p_{v_1},\dots,p_{v_n})$.
    \ii If $v,Tv,\dots,T^{n-1}v$ is a basis ($v$ is a \vocab{cyclic vector}), then $m_T=p_v$, which has degree $n$.
  \end{enumerate}
\end{proposition*}
\begin{ideabox}
To show: $m_T$ from computations with single vectors, which are cheaper than powers of a matrix.

Polynomials that kill one vector behave like polynomials that kill $T$: the division argument of Prop.~65 works the same way. An operator that is $0$ on a basis is $0$.

Killing $T$ means killing every basis vector, i.e.\ being divisible by every $p_{v_i}$. The least such polynomial is the lcm.
\end{ideabox}
\begin{proof}
  \begin{steps}
    \ii (a) Divide: $q=p_vs+r$ with $\deg r<\deg p_v$ or $r=0$. If $q(T)v=0$, then $r(T)v=q(T)v-s(T)p_v(T)v=0$, so $r=0$ by minimality of $p_v$. Since $m_T(T)v=0$, $p_v\mid m_T$.
    \ii (b) Let $L$ be the lcm. Each $p_{v_i}\mid L$, so $L(T)v_i=0$ by (a); $L(T)$ is $0$ on a basis, so $L(T)=0$ and $m_T\mid L$.
    \ii Each $p_{v_i}\mid m_T$ by (a), so $L\mid m_T$. Both are monic, so $L=m_T$.
    \ii (c) $v,\dots,T^{n-1}v$ are independent, so no nonzero polynomial of degree $<n$ kills $v$: $\deg p_v\ge n$. Also $p_v\mid m_T$ and $\deg m_T\le n$ (Prop.~67). So $m_T=p_v$.
  \end{steps}
\end{proof}

\begin{proposition*}[F8: real versus complex]
  For a real matrix $A$, $m_A$ is the same over $\RR$ and over $\CC$. So $A$ is diagonalizable over $\CC$ $\iff m_A$ has distinct roots, and diagonalizable over $\RR$ $\iff$ in addition every root is real.
\end{proposition*}
\begin{ideabox}
To show: Computing over $\CC$ cannot give a different answer than over $\RR$.

Def.~36 finds $m_A$ by solving linear equations whose coefficients are entries of powers of $A$, all real. Solving linear equations never leaves the real numbers.
\end{ideabox}
\begin{proof}
  \begin{steps}
    \ii For each $k$, ``$A^k=a_{k-1}A^{k-1}+\dots+a_0I$'' is a system of $n^2$ linear equations in $a_0,\dots,a_{k-1}$ with real coefficients.
    \ii Row reduction uses only $+,-,\times,\div$ of real numbers, so the system has a complex solution iff it has a real one.
    \ii So the first solvable $k$ is the same over $\RR$ and $\CC$, and the solution is unique (Def.~36), hence real. So $m_A$ is the same.
    \ii The consequences are Prop.~68 over each field: over $\RR$ the linear factors must have real roots.
  \end{steps}
\end{proof}

\begin{proposition*}[F9: projections onto eigenspaces, HW7 P7]
  If $m_T=\prod_{i=1}^m(x-\lambda_i)$ with distinct roots, then $g_j(T)$, with $g_j$ as in Prop.~68, is the projection onto $E(\lambda_j,T)$ along $\bigoplus_{i\ne j}E(\lambda_i,T)$.
\end{proposition*}
\begin{proof}
  On $E(\lambda_i,T)$, $g_j(T)$ acts as the number $g_j(\lambda_i)$ (T7), which is $1$ for $i=j$ and $0$ for $i\ne j$. Since $V=\bigoplus_iE(\lambda_i,T)$, $g_j(T)(v_1+\dots+v_m)=v_j$.
\end{proof}

\begin{proposition*}[F10: repeated roots, Ex.~59]
  \begin{enumerate}[(a)]
    \ii $\lambda$ is a repeated root of $p$ $\iff p(\lambda)=p'(\lambda)=0$.
    \ii The common roots of $p$ and $p'$ are the roots of $\gcd(p,p')$, which the Euclidean algorithm finds without factoring $p$.
    \ii Over $\CC$, if $T^k=I$ for some $k\ge1$, then $T$ is diagonalizable.
  \end{enumerate}
\end{proposition*}
\begin{ideabox}
To show: A test for repeated roots that does not need the roots.

From calculus: at a double root the graph touches the axis, so the value and the slope are both $0$.

The product rule turns this into algebra: differentiating $(x-\lambda)^2q$ leaves a factor $x-\lambda$.
\end{ideabox}
\begin{proof}
  (a)
  \begin{steps}
    \ii If $p=(x-\lambda)^2q$, then $p'=2(x-\lambda)q+(x-\lambda)^2q'$, which is $0$ at $\lambda$.
    \ii Conversely, if $p(\lambda)=0$, write $p=(x-\lambda)q$ (T6). Then $p'=q+(x-\lambda)q'$, so $p'(\lambda)=q(\lambda)$. If this is $0$, then $(x-\lambda)\mid q$, so $(x-\lambda)^2\mid p$.
  \end{steps}
  (b) Divide $p$ by $p'$, then $p'$ by the remainder, and so on. Each remainder is (previous dividend) $-$ (quotient)(previous divisor), so a common root of $p,p'$ is a root of every remainder, in particular of the last nonzero one, $g=\gcd(p,p')$. Conversely, $g$ divides $p$ and $p'$ (work back up the divisions).

  Example: $p=x^3-3x+2$, $p'=3x^2-3$. Then $p=\frac x3\,p'+(-2x+2)$, and $p'=-\frac32(x+1)\cdot(-2x+2)$ with remainder $0$. So $g$ is a multiple of $x-1$, and $1$ is a double root: $p=(x-1)^2(x+2)$.

  (c)
  \begin{steps}
    \ii $T^k=I$ means $x^k-1$ kills $T$, so $m_T\mid x^k-1$.
    \ii $(x^k-1)'=kx^{k-1}$ has only the root $0$, which is not a root of $x^k-1$. By (a), $x^k-1$ has no repeated root, so it is a product of $k$ distinct linear factors.
    \ii $m_T$ is a monic divisor of it, so it is a product of distinct linear factors (F3(b)). Prop.~68.
  \end{steps}
\end{proof}

\begin{proposition*}[F11: odd dimension, LADR 5B]
  Every operator on an odd-dimensional real vector space has an eigenvalue.
\end{proposition*}
\begin{ideabox}
To show: A real eigenvalue, i.e.\ a real root of $m_T$ (Prop.~66).

Over $\RR$ there may be none: the rotation of $\RR^2$ has $m=x^2+1$. But $\RR^2$ has even dimension.

Plan: Show that ``no real root'' forces the dimension to be even, in three steps. (A) A real polynomial with no real root is a product of quadratics $x^2+bx+c$ with $b^2<4c$. (B) For such a quadratic $p$, $\dim\nul p(T)$ is even: each nonzero $u$ in it comes with $Tu$, and $\spn(u,Tu)$ is $2$-dimensional. (C) Split off $\nul p(T)$ and use induction on the dimension.
\end{ideabox}
\needspace{7\baselineskip}\textbf{Step A: real polynomials.} A real polynomial is a product of real linear factors and real quadratics $x^2+bx+c$ with $b^2<4c$.
\begin{steps}
  \ii For real $p$, $p(\bar z)=\overline{p(z)}$, so if $z$ is a non-real root, so is $\bar z$.
  \ii $q=(x-z)(x-\bar z)=x^2-2(\operatorname{Re}z)x+|z|^2$ is real, and $b^2-4c=-4(\operatorname{Im}z)^2<0$.
  \ii Divide over $\RR$: $p=qs+r$ with $\deg r<2$. Then $r(z)=r(\bar z)=0$, and a polynomial of degree $<2$ with two roots is $0$. So $p=qs$ with $s$ real. Real roots split off real linear factors (T6). Repeat on $s$ (induction on the degree).
\end{steps}
\needspace{7\baselineskip}\textbf{Step B: if $b^2<4c$ and $p=x^2+bx+c$, then $\dim\nul p(T)$ is even.} Induction on $\dim\nul p(T)$; let $U=\nul p(T)$, which is invariant (T5).
\begin{steps}
  \ii $T|_U$ has no eigenvalue: $Tu=\lambda u$ with $u\in U$, $u\ne0$, gives $0=p(T)u=p(\lambda)u$, but $p(\lambda)=(\lambda+\frac b2)^2+c-\frac{b^2}4>0$.
  \ii If $U=0$, its dimension $0$ is even. Otherwise pick $u\in U$, $u\ne0$, and let $W=\spn(u,Tu)$. It has dimension $2$, since $u$ is not an eigenvector. It is invariant, since $T(Tu)=-bTu-cu$.
  \ii The quotient operator $S=T|_U/W$ on $U/W$ satisfies $p(S)=0$ (T10), so $\nul p(S)$ is all of $U/W$, of dimension $\dim U-2$. By induction this is even, so $\dim U$ is even.
\end{steps}
\needspace{7\baselineskip}\textbf{Step C: proof of F11.} Induction on the odd number $n=\dim V$; $n=1$ is clear.
\begin{steps}
  \ii Suppose $T$ has no eigenvalue. Then $m_T$ has no real root (Prop.~66), so by Step A, $m_T=ps$ with $p=x^2+bx+c$, $b^2<4c$.
  \ii $s(T)\ne0$ (smaller degree than $m_T$), so some $w$ has $s(T)w\ne0$, and $p(T)\bigl(s(T)w\bigr)=0$. So $\nul p(T)\ne0$, and its dimension is even by Step B.
  \ii $W=\range p(T)$ is invariant, and $\dim W=n-\dim\nul p(T)$ is odd and less than $n$ (T1). $T|_W$ has no eigenvalue, because an eigenvector of $T|_W$ is one of $T$. This contradicts the induction hypothesis.
\end{steps}


\subsubsection*{Computing $m_T$}
Over $\CC$, or whenever $m_T$ splits into linear factors, $m_T=\prod_i(x-\lambda_i)^{k_i}$, where $\lambda_1,\dots,\lambda_m$ are the distinct eigenvalues (Prop.~66), every $k_i\ge1$, and $\sum_ik_i\le n$ (Prop.~67). So the work is finding the eigenvalues and then the exponents. Use whichever method below needs the least computation.

\method{Method 1: from the definition.} Compute $I,A,A^2,\dots$ and stop at the first power that is a combination of the earlier ones. Best for $2\times2$ matrices or when powers are easy.
\begin{example*}
  $A=\begin{psmallmatrix}2&1\\0&3\end{psmallmatrix}$ is not a multiple of $I$, so $I,A$ are independent. $A^2=\begin{psmallmatrix}4&5\\0&9\end{psmallmatrix}$. Solve $A^2=aA+bI$ entry by entry: entry $(1,2)$ gives $a=5$; entry $(1,1)$ gives $4=2a+b$, so $b=-6$; entry $(2,2)$ checks: $9=3a+b$. So $m_A=x^2-5x+6=(x-2)(x-3)$.
\end{example*}

\method{Method 2: eigenvalues, then test candidates.} Find the eigenvalues (for a triangular matrix, the diagonal entries, Prop.~60). The candidates are $\prod_i(x-\lambda_i)^{k_i}$ with every $k_i\ge1$ and $\sum k_i\le n$. Test them in increasing degree; the first one that kills $A$ is $m_A$. Reason: $m_A$ is one of the candidates, and any candidate $q$ with $q(A)=0$ is a multiple of $m_A$ (Prop.~65), so $\deg q>\deg m_A$ unless $q=m_A$.
\begin{example*}
  $A=\begin{psmallmatrix}3&1&1\\0&3&1\\0&0&1\end{psmallmatrix}$ has eigenvalues $3,3,1$. Candidates: degree $2$: $(x-3)(x-1)$; degree $3$: $(x-3)^2(x-1)$ and $(x-3)(x-1)^2$. $(A-3I)(A-I)=\begin{psmallmatrix}0&2&1\\0&0&0\\0&0&0\end{psmallmatrix}\ne0$, and $(A-3I)^2(A-I)=0$. So $m_A=(x-3)^2(x-1)$, and $A$ is not diagonalizable.

  $B=\begin{psmallmatrix}1&0&1\\0&1&0\\0&0&2\end{psmallmatrix}$ has eigenvalues $1,1,2$, and $(B-I)(B-2I)=0$. So $m_B=(x-1)(x-2)$, and $B$ is diagonalizable although it is not diagonal.
\end{example*}

\method{Method 3: blocks.} For a block diagonal matrix, $m$ is the lcm of the blocks' minimal polynomials (F4). For a block upper-triangular matrix $A=\begin{psmallmatrix}B&*\\0&C\end{psmallmatrix}$: $m_B\mid m_A$, $m_C\mid m_A$ and $m_A\mid m_Bm_C$. So if $m_B$ and $m_C$ have no common root, $m_A=m_Bm_C$.
\why{\begin{steps}
  \ii The first block of basis vectors spans an invariant subspace $U$ with $B=\MM(T|_U)$. $C$ is the matrix of $T/U$ in the cosets of the remaining basis vectors.
  \ii $m_B\mid m_A$ and $m_C\mid m_A$ by F3(a). $m_A\mid m_Bm_C$ by step 4 of the proof of Prop.~67, which shows that $m_{T|_U}m_{T/U}$ kills $T$.
  \ii ``No common root'' means no common root in $\CC$. Over $\CC$, $m_A=\prod(x-\lambda)^{e_\lambda}$, and each monic divisor of it is $\prod(x-\lambda)^{b_\lambda}$ with $b_\lambda\le e_\lambda$ (its roots are roots of $m_A$; divide out one factor at a time, T6).
  \ii So $m_B$ and $m_C$ are such products over disjoint sets of $\lambda$'s, and their product divides $m_A$.
\end{steps}}
\begin{itemize}
  \ii $A=\begin{psmallmatrix}1&-2&0&0\\1&-1&0&0\\0&0&2&0\\0&0&0&2\end{psmallmatrix}$: the first block $R$ has $R^2=-I$ and is not a multiple of $I$, so $m_R=x^2+1$; the block $2I$ has $m=x-2$. So $m_A=(x-2)(x^2+1)$. Over $\CC$ the roots $2,\pm i$ are distinct, so $A$ is diagonalizable over $\CC$; over $\RR$ it is not, since $\pm i\notin\RR$ (F8).
  \ii The matrix $A$ of Method 2 is $\begin{psmallmatrix}B&*\\0&1\end{psmallmatrix}$ with $B=\begin{psmallmatrix}3&1\\0&3\end{psmallmatrix}$, $m_B=(x-3)^2$. The roots $3$ and $1$ differ, so $m_A=(x-3)^2(x-1)$ with no matrix multiplication.
\end{itemize}

\method{Method 4: follow one vector (F7).} Compute $v,Av,A^2v,\dots$ until the first dependence; it gives $p_v$, which divides $m_A$. If $\deg p_v=n$, then $m_A=p_v$. Otherwise repeat with more basis vectors and take the lcm.
\begin{example*}
  $A=\begin{psmallmatrix}0&0&2\\1&0&-1\\0&1&2\end{psmallmatrix}$: $Ae_1=e_2$, $Ae_2=e_3$, so $e_1,Ae_1,A^2e_1$ is a basis. $A^3e_1=Ae_3=(2,-1,2)=2e_1-Ae_1+2A^2e_1$. So $m_A=x^3-2x^2+x-2=(x-2)(x^2+1)$.

  $A=\begin{psmallmatrix}2&0&0\\1&2&0\\0&0&3\end{psmallmatrix}$: $Ae_1=(2,1,0)$, $A^2e_1=(4,4,0)=4Ae_1-4e_1$, so $p_{e_1}=(x-2)^2$. $Ae_3=3e_3$, so $p_{e_3}=x-3$. The lcm $(x-2)^2(x-3)$ has degree $3=n$, so it is $m_A$; there is no need to compute $p_{e_2}$.
\end{example*}
The first matrix is a \vocab{companion matrix}: if $Ae_j=e_{j+1}$ for $j<n$ and $Ae_n=a_0e_1+\dots+a_{n-1}e_n$, then $e_1$ is cyclic and $A^ne_1=Ae_n$, so $m_A=x^n-a_{n-1}x^{n-1}-\dots-a_0$. So every monic polynomial is the minimal polynomial of some matrix.

\method{Method 5: spot a relation.} If you notice $p(A)=0$, then $m_A\mid p$; check which divisors work.
\begin{itemize}
  \ii Rank one, $A=uv^{\mathsf T}$ (all columns multiples of $u$): $A^2=u(v^{\mathsf T}u)v^{\mathsf T}=(\tr A)A$, where $\tr A$ is the sum of the diagonal entries, which equals $v^{\mathsf T}u$. For $n\ge2$, $A$ is not a multiple of $I$, so $m_A=x(x-\tr A)$ if $\tr A\ne0$ (diagonalizable) and $m_A=x^2$ if $\tr A=0$ (not diagonalizable). Example: $\begin{psmallmatrix}1&-1\\-1&1\end{psmallmatrix}^2=2\begin{psmallmatrix}1&-1\\-1&1\end{psmallmatrix}$, so $m=x(x-2)$.
  \ii $A^2=A$: $m_A$ is $x$, $x-1$ or $x(x-1)$. $A^2=I$: $m_A$ is $x-1$, $x+1$ or $(x-1)(x+1)$. Both cases are diagonalizable. $A^k=0$: $m_A=x^j$ for the least $j$ with $A^j=0$.
  \ii A permutation matrix satisfies $P^k=I$ for some $k\ge1$ (only finitely many permutations, so two powers agree), so it is diagonalizable over $\CC$ (F10).
\end{itemize}

\method{Method 6: the $2\times2$ formula.} For $2\times2$ $A=\begin{psmallmatrix}a&b\\c&d\end{psmallmatrix}$, with $\tr A=a+d$ and $\det A=ad-bc$: $m_A=x^2-(\tr A)x+\det A$, unless $A=cI$ (then $m_A=x-c$).
\begin{itemize}
  \ii Why: for $A=\begin{psmallmatrix}a&b\\c&d\end{psmallmatrix}$, $A^2-(a+d)A=\begin{psmallmatrix}bc-ad&0\\0&bc-ad\end{psmallmatrix}$, so $A^2-(\tr A)A+(\det A)I=0$. If $A\ne cI$, $\deg m_A=2$ and $m_A$ divides this monic quadratic, so they are equal.
  \ii Example: $A=\begin{psmallmatrix}1&4\\-1&-3\end{psmallmatrix}$: $\tr A=-2$, $\det A=1$, so $m_A=(x+1)^2$; $A\ne-I$, so $A$ is not diagonalizable.
\end{itemize}

\method{Method 7: operators on function spaces.} Work with the operator directly, or take its matrix in a basis.
\begin{example*}
  $D$ on $V=\spn(e^x,xe^x)$: $(D-I)e^x=0$, $(D-I)(xe^x)=e^x\ne0$, and $(D-I)^2(xe^x)=0$. So $(D-I)^2=0$ and $D\ne I$: $m_D=(x-1)^2$, and $D$ is not diagonalizable.
\end{example*}

\begin{shortcuts}
  \ii \textbf{No computation needed.} $n$ distinct eigenvalues (e.g.\ a triangular matrix with distinct diagonal entries) $\Rightarrow m=\prod(x-\lambda_i)$ (every eigenvalue is a root and $\deg m\le n$), diagonalizable. One eigenvalue $\lambda$ and $A\ne\lambda I$ $\Rightarrow$ not diagonalizable.
  \ii \textbf{Testing a candidate}: one nonzero entry is enough to show it fails.
  \ii \textbf{Inverse} (F1). $A=\begin{psmallmatrix}2&1\\0&3\end{psmallmatrix}$: $A^2-5A+6I=0$, so $A(5I-A)=6I$ and $A^{-1}=\frac16(5I-A)$.
  \ii \textbf{High powers.} Divide: $x^k=m(x)q(x)+r(x)$ with $\deg r<\deg m$; then $A^k=r(A)$. Find $r$ by putting each root $\lambda$ of $m$ into this identity: $r(\lambda)=\lambda^k$. For a double root, differentiate as well: $r'(\lambda)=k\lambda^{k-1}$ (the derivative of $(x-\lambda)^2q$ is $2(x-\lambda)q+(x-\lambda)^2q'$, which vanishes at $\lambda$).
  \begin{itemize}
    \ii $A=\begin{psmallmatrix}2&1\\0&3\end{psmallmatrix}$, $r=ax+b$: $2a+b=2^k$, $3a+b=3^k$, so $A^k=(3^k-2^k)A+(3\cdot2^k-2\cdot3^k)I$.
    \ii $m=(x+1)^2$ (as for $\begin{psmallmatrix}1&4\\-1&-3\end{psmallmatrix}$): $r(-1)=(-1)^k$, $r'(-1)=k(-1)^{k-1}$, so $A^k=(-1)^kI+k(-1)^{k-1}(A+I)$.
  \end{itemize}
  \ii \textbf{Any polynomial} $p(A)$ equals $r(A)$, where $r$ is the remainder of $p$ divided by $m_A$. So every polynomial in $A$ is a combination of $I,A,\dots,A^{d-1}$.
  \ii \textbf{Eigenvectors without solving} (Prop.~66). If $m=(x-\lambda)s$, the nonzero columns of $s(A)$ are eigenvectors for $\lambda$; if $A$ is diagonalizable, they span $E(\lambda)$ (on $E(\mu)$, $s(A)$ is the number $s(\mu)$, which is $0$ for $\mu\ne\lambda$ and nonzero for $\mu=\lambda$). For $A=\begin{psmallmatrix}2&1\\0&3\end{psmallmatrix}$ and $\lambda=3$: $s(A)=A-2I=\begin{psmallmatrix}0&1\\0&1\end{psmallmatrix}$, so $(1,1)$ is an eigenvector for $3$.
  \ii \textbf{Diagonalizable without finding eigenvalues.} One polynomial that kills $T$ and has distinct roots in $\FF$ is enough, since $m_T$ divides it: $T^2=T$, $T^2=I$, $T^3=T$ (over $\RR$), $T^k=I$ (over $\CC$).
  \ii \textbf{Not similar}: different minimal polynomials prove it (F5). Equal minimal polynomials do not prove similarity.
  \ii \textbf{Lower bound}: if $v,Av,\dots,A^{j-1}v$ are independent, then $\deg m_A\ge j$ (F7).
\end{shortcuts}

\begin{trapbox}
  \ii Only roots in $\FF$ are eigenvalues: over $\RR$, $m=x^2+1$ gives no eigenvalues.
  \ii A polynomial that kills $T$ need not be $m_T$: $T^2=T$ allows $m_T=x$, $x-1$ or $x^2-x$. Check the divisors.
  \ii Every eigenvalue is a root of $m_T$, so a candidate that leaves out an eigenvalue never works.
  \ii $p(T)v=0$ for one $v\ne0$ only gives $p_v\mid p$, not $m_T\mid p$. And $v$ need not be an eigenvector even when $p$ has distinct roots: $T=\operatorname{diag}(1,2)$, $v=(1,1)$, $p=(x-1)(x-2)$. What does follow, if $p$ splits into linear factors over $\FF$ (always true over $\CC$), is that some root of $p$ is an eigenvalue (T4, T6). Over $\RR$ not even that: $T=\begin{psmallmatrix}0&-1\\1&0\end{psmallmatrix}$, $p=x^2+1$, $v=e_1$.
  \ii Over $\RR$, $T^k=I$ does not imply diagonalizable: $R=\begin{psmallmatrix}1&-2\\1&-1\end{psmallmatrix}$ has $R^4=I$ but $m_R=x^2+1$.
  \ii $m_{T^{-1}}(x)$ is not $m_T(1/x)$, which is not a polynomial. Reverse the coefficients and divide by $c_0$.
  \ii $\deg m_T=n$ says nothing about diagonalizability: $\operatorname{diag}(1,2)$ and $J_2(0)$ both have degree $2$.
\end{trapbox}



\section{Problem set}

\begin{prob}
\textbf{HW7 P1. }\ Return to the matrices from Homework 6, $A=\begin{pmatrix}0&-1\\1&0\end{pmatrix}$ and $B=\begin{pmatrix}1&1\\0&1\end{pmatrix}$.
\begin{enumerate}[(a)]
  \ii Find the minimal polynomial of $A$ over $\RR$ and over $\CC$.
  \ii Find the minimal polynomial of $B$ over $\RR$.
  \ii Use the minimal-polynomial criterion for diagonalizability to recover the conclusions of Homework 6 about whether these operators are diagonalizable over the indicated fields.
\end{enumerate}
\end{prob}

\begin{prob}
\textbf{HW7 P2. }\ Let $D\colon\PP_3(\RR)\to\PP_3(\RR)$ be differentiation.
\begin{enumerate}[(a)]
  \ii Find the minimal polynomial of $D$.
  \ii Decide whether $D$ is diagonalizable.
\end{enumerate}
\end{prob}

\begin{prob}
\textbf{HW7 P3. }\ Let $T$ be an invertible operator on a finite-dimensional vector space, and let $m_T(x)=x^d+a_{d-1}x^{d-1}+\dots+a_1x+a_0$ be its minimal polynomial.
\begin{enumerate}[(a)]
  \ii Prove that $a_0\ne0$.
  \ii Prove that $T^{-1}$ is a polynomial in $T$.
  \ii Suppose now that $V$ is a complex vector space and that $m_T(x)=x^4+x^3-x^2-5x+4$. Express $T^{-1}$ explicitly as a polynomial in $T$, find the eigenvalues of $T$, and decide whether $T$ is diagonalizable.
\end{enumerate}
\end{prob}

\begin{prob}
\textbf{HW7 P4. (optional)}\ Let $m_T$ be the minimal polynomial of an operator $T$.
\begin{enumerate}[(a)]
  \ii Prove that the minimal polynomial of $T-\lambda I$ is $m_T(x+\lambda)$.
  \ii If $c\ne0$ and $d=\deg m_T$, prove that the minimal polynomial of $cT$ is $c^dm_T(x/c)$.
\end{enumerate}
\end{prob}

\begin{prob}
\textbf{HW7 P5. }\ Let $A$ be a $2\times2$ complex matrix whose only eigenvalue is $\lambda$. Prove that its minimal polynomial is either $x-\lambda$ or $(x-\lambda)^2$. Show that the first case occurs exactly when $A=\lambda I$, and conclude that $A$ is diagonalizable if and only if $A=\lambda I$.
\end{prob}

\begin{prob}
\textbf{HW7 P6. }\ Consider $C=\begin{pmatrix}2&1&0\\0&2&0\\0&0&3\end{pmatrix}$. Find its eigenvalues, eigenspaces, and minimal polynomial. Decide whether $C$ is diagonalizable.
\end{prob}

\begin{prob}
\textbf{HW7 P7. (optional)}\ Let $T$ be an operator whose minimal polynomial splits as $m_T(x)=(x-\lambda_1)\cdots(x-\lambda_r)$ with distinct roots. For a fixed $j$, construct a polynomial $p_j$ such that $p_j(\lambda_j)=1$ and $p_j(\lambda_k)=0$ for $k\ne j$, and explain why $p_j(T)$ is the projection onto $E(\lambda_j,T)$ along the sum of the other eigenspaces.
\end{prob}

\begin{prob}
\textbf{M1. Computing minimal polynomials.}\ Find the minimal polynomial and decide whether the matrix is diagonalizable.
\begin{enumerate}[(a)]
  \ii $A=\begin{pmatrix}1&1&1\\1&1&1\\1&1&1\end{pmatrix}$.
  \ii $P=\begin{pmatrix}0&0&1\\1&0&0\\0&1&0\end{pmatrix}$, over $\RR$ and over $\CC$.
  \ii $B=\begin{pmatrix}3&1&0&0\\0&3&0&0\\0&0&3&0\\0&0&0&5\end{pmatrix}$.
\end{enumerate}
\end{prob}

\begin{prob}
\textbf{M2. Operators on polynomials.}\ Let $n\ge1$.
\begin{enumerate}[(a)]
  \ii Let $T\colon\PP_n(\RR)\to\PP_n(\RR)$, $T(p)=xp'$. Find $m_T$ and decide whether $T$ is diagonalizable.
  \ii Let $S\colon\PP_n(\RR)\to\PP_n(\RR)$, $(Sp)(x)=p(x+1)$. Prove that $m_S=(x-1)^{n+1}$, and decide whether $S$ is diagonalizable.
\end{enumerate}
\end{prob}

\begin{prob}
\textbf{M3. }\ Let $V$ be a finite-dimensional real vector space and $T\in\LL(V)$ with $T^3=T$. Prove that $T$ is diagonalizable and that every eigenvalue of $T$ is $-1$, $0$ or $1$.
\end{prob}

\begin{prob}
\textbf{M4. }\ Let $V\ne\{0\}$ be finite-dimensional and $T\in\LL(V)$. Prove that $\deg m_T=1$ if and only if $T=cI$ for some $c\in\FF$.
\end{prob}

\begin{prob}
\textbf{M5. Minimal polynomial of the inverse.}\ Let $T$ be invertible with $m_T(x)=x^d+a_{d-1}x^{d-1}+\dots+a_0$. Prove that $m_{T^{-1}}(x)=\frac1{a_0}x^dm_T(1/x)$.
\end{prob}

\begin{prob}
\textbf{M9. Using the minimal polynomial to compute.}\ Let $A=\begin{pmatrix}1&2\\3&4\end{pmatrix}$. Find $m_A$, write $A^{-1}$ as a polynomial in $A$, and write $A^3$ as $aA+bI$.
\end{prob}

\end{document}
